\section{Preliminaries}
\subsection{Adaptive Evolving Neural Security Architectures}
Adaptive evolving neural security architectures refer to intelligent systems that leverage neural networks capable of learning and self-improvement over time. These architectures continuously analyze new data and threat patterns, allowing them to autonomously adjust their defensive strategies. This dynamic approach is particularly effective in environments where cyber threats are constantly changing, such as cloud and IoT ecosystems.

\subsection{Cloud-Native CI/CD Security and Real-Time Defense}
Cloud-native CI/CD security integrates protective measures directly into the software development and deployment pipeline. By embedding security checks and automated defenses within continuous integration and continuous deployment (CI/CD) processes, organizations can identify and address vulnerabilities as code is developed, tested, and released. Real-time defense systems help keep applications safe while they’re running, making sure security keeps up with the fast and scalable nature of cloud environments.

\subsection{Polymorphic and Metamorphic Malware Detection}
Polymorphic malware keeps its payload the same but changes its external form, while Metamorphic malware rewrites its own code through each generation to change the payload. Detection techniques that focus on behavioral analysis and self-learning techniques can recognize these threats better than solely relying on static signatures

\subsection{Behavioral and Code-Level Anomaly Detection}
Behavioral detection techniques involve differentiating the normal runtime behaviour and activities that deviate from the learned baseline from normal runtimes, while Code-level anomaly detection inspects the structure and instructions of the code to spot abnormalities.


\subsection{Static and Dynamic Threat Analysis (SAST/DAST)}
SAST involves analyzing without running the app to find vulnerabilities in code, bytecode, or binaries, while DAST probes a running application from outside to find exploitable flaws.

\subsection{Adversarial Training}
Adversarial, in a general sense, means introducing an opposition that tries to fool the system. Adversarial training is a method where models are trained by inputting intentionally perturbed samples to make the model predict wrong and learn decision boundaries 


\subsection{Red Team Simulation}
Red team simulation is a security exercise where authorized attackers emulate real adversaries to test how well a system performs its security procedures and its detection technology.


\section{Review of Existing Research}
To systematically organize the literature, twenty research papers were reviewed and grouped into five key themes based on their focus, methods, and relevance to the proposed adaptive threat detection framework. This structure enables clear comparison across studies that address different aspects of cloud CI/CD security and intelligent defense. The themes include: (1) Adaptive AI and Neural Architecture Design, exploring self-evolving neural models; (2) Cloud and CI/CD Pipeline Security, centered on protecting DevOps workflows and automated releases; (3) Threat Characterization and Detection, analyzing malware behaviors and detection methods; (4) Analytical and Testing Methodologies, focused on SAST and DAST-driven vulnerability assessments; and (5) Adversarial and Simulation-Based Defense, investigating Red Teaming and adversarial learning techniques. Organizing the research in this way helps reveal how each area contributes to AI-driven cloud security and exposes the critical gaps this study seeks to address.\\

\subsection{Adaptive AI and Neural Architecture Design}
The study aims to highlight the advantages of ENN, Evolving Neural Networks, for decision-making for effective discovery and adaptive behaviors across complex domains. They employed Neuroevolution techniques to optimize network parameters and architectures simultaneously but separately. They used CNE focusing on evolving connection weights of fixed topology networks through genetic evolution and Advanced Neuroevolution for individual neurons, partial networks, and evolving network topology simultaneously with NEAT. They also used Indirect Encoding to efficiently represent large networks with reusable structure, allowing scaling to millions of connections. The combined Neuroevolution with supervised learning algorithms to fine-tune network weights. Through their combined system, they found that it outperforms traditional reinforced learning approaches by a magnitude in complex control problems. They found that Neuevolution demonstrates strong adaptability and robustness in dynamic and non-Markovian environments, especially when combined with supervised learning, which further enhances adaptability(\textbf{Miikkulainen \& Stanley, 2009})\cite{Miikkulainen}.\\

The research focuses on evolving ANNs for adaptive behavior in artificial agents, arguing that evolving both the architecture and learning parameters incrementally allows for the emergence of a robust, noise-resistant network that is capable of complex, adaptive behaviors in dynamic environments. This study models recurrent neural networks with excitation and inhibition channels, nonlinear transfer functions, veto mechanisms, and internal noise to emulate biological neural variability to enhance adaptability. They use an extended genetic algorithm, SA-GA, allowing both connectivity and network size to grow as needed. They start with minimal random changes in connected network populations, increasing complexity only when beneficial to fitness, measured by a robot’s ability to navigate to and stay centered in an arena using limited sensors. The metrics and analysis reveal that this evolutionary process promotes modularity and the emergence of a biologically inspired adaptive control system(\textbf{Harvey \& Husbands, 1992})\cite{Harvey}.\\

Another study explores the use of Neural Architecture Search (NAS) to discover inherently robust neural network architectures that can resist adversarial attacks. This research proposes an evolutionary NAS approach, Robust Neural Architecture Search (R-NAS), which simultaneously optimizes for accuracy and robustness against adversarial attacks using a combined fitness function. The neural architectures are evolved across layers, blocks, and model level by diverse mutation operations and spectrum-based niching to maintain population diversity. By testing each design against black-box attacks on the CIFAR-10 dataset, it balances high performance on normal data with strong resistance to adversarial interference. R-NAS created neural networks that stayed highly accurate while being more than twice as resistant to adversarial attacks compared to traditional methods(\textbf{Kotyan \& Vargas, 2020})\cite{Kotyan}.\\

This investigation introduces an adaptive neural network model with reflection capable of autonomously evolving its internal structure and learning parameters in response to dynamic environments, modeling how biological neural networks evolve and adapt both their architecture and behavior, overcoming the limitations of fixed-topology neural networks. It introduces neuro-agents, modular units made up of smaller components called neuro-objects that learn together by sharing information and working in parallel. These modules can be organized hierarchically, combining similar and diverse learning parts to form large, flexible networks that adapt and work together seamlessly on complex problems. The study shows that neural networks can grow and adapt both their design and learning at once, combining specialized parts into large systems that work together without losing their ability to learn(\textbf{Tsukamoto \& Namatame, 1996})\cite{Tsukamoto}.\\

\subsection{Cloud and CI/CD Pipeline Security}
This paper used the CICD2019 dataset to detect threats. It encoded the CICD logs to an Autoencoder (AE) so that it can feed insightful features of this dataset to a Multi-layer Perception (MLP) Deep Learning model that can detect threats. Then the threats were categorized into five categories.\\

\textbf{Normal}: expected network behavior.\\
\textbf{DoS (Denial of Service)}: Floods the network with traffic\\
\textbf{Probe}: Scanning aimed at discovering vulnerabilities.\\
\textbf{U2R (User to Root)}: Attempt to gain access to root by using a privileged user like admin.\\
\textbf{R2L (Remote to Local)}: Remotely making unauthorized attempts to gain local network access.\\
 (\textbf{Modalavalasa \& Yadav, 2025})\cite{Modalavalasa}\\
 
This paper uses Docker image scanning tools to detect malicious activity. Uses Snyk and StackHawk to detect malicious activity throughout the application development cycle. This paper combines two technologies to improve CICD security. They proposed both SAST and DAST type testing for more secure pipeline security.\\
\begin{enumerate}
    \item Snyk's focus on SAST (Static Application Security Testing)
    \item StackHawk's DAST (Dynamic Application Security Testing)
\end{enumerate}
(\textbf{Marandi, Bertia, \& Silas, 2023})\cite{Marandi}\\

In this paper, when pushing code in GitHub, it automatically triggered Jenkins, then it fetched the code, ran OWASP Dependency Checker, SonarQube, and Trivy to ensure code quality and check vulnerabilities. Then it was pushed to a Docker image and deployed through AgroCD. AgroCD checks configuration requirements before deployment. AWS VPC isolates sensitive data in different subnets to avoid public exposure. Monitoring was done by Prometheus and CloudWatch. They provided real-time detection of security issues across the pipeline (\textbf{Ashtagi et al., 2025})\cite{Ashtagi}.\\

This paper broadly focuses on how software development can be modernized with the integration of governance and risk management. From development cycles to deployments and integrations, it proposes a unified framework combining SDLC, CI/CD, and DevOps. For security, it includes Security methods, including SAST, DAST, SCA, Threat Modeling, Risk-Based Testing, and penetration testing (\textbf{Ganugapati \& Prabhu, 2023})\cite{Ganugapati}.\\

\subsection{Threat Characterization and Detection}
The research introduces a novel malware detection system that addresses the shortcomings associated with existing signature-based scanners, which are susceptible to typical obfuscation techniques such as code reordering and register renaming. It uses instruction sequence templates to represent malware behaviors and employs an intermediate representation of machine instructions for architecture-independent analysis. A semantics-aware detecting algorithm performs symbolic execution and theorem proving for verification of semantic equivalence while avoiding def-use chains corruption. The results of the experiment demonstrate that using this approach, it is possible to detect several worm variants using far fewer signatures and without any false positives on benign software, thus substantially surpassing standard static analyzers when it comes to tackling obfuscated code. Nevertheless, the proposed approach has some limitations: first, they should be successfully disassembled in order to analyze their instructions; second, there are supported instruction forms; and lastly, there are strict matching rules. This may cause a mis-catch for some variations. However, all these factors lead us to concentrate on Windows IA-32 binaries (\textbf{Christodorescu, Jha, Seshia, Song, \& Bryant, 2005})\cite{Christodorescu}.\\

This study introduces a deep learning method that combines convolutional and recurrent neural networks to classify malware on the basis of system call sequences, which are collected at runtime in a sandbox. A simple list of API calls can be converted to their corresponding scalar representations; it reveals the hierarchical structure of API calls in software and improves the accuracy for discriminating malware families. When tested on a balanced dataset created from antivirus signatures, the ConvNet-LSTM model showed over 89\% accuracy, which is a better performance than other machine learning models. Nonetheless, this model only considers system calls, may have imprecise labels and suffers from possible high computational requirements of dynamic analysis as well as evasion attacks (\textbf{Kolosnjaji, Zarras, Webster, \& Eckert, 2016})\cite{Kolosnjaji}.\\

Another investigation holds the core information about the methods that have been developed by malware over a period of 10 years to evade dynamic tools like sandboxes and debugging software. Based on this, it enlists 92 different techniques in 16 groups and surveys more than 45k malware samples from the time frame of 2010–2019 against benign software. The results demonstrate that approximately 80\% of malware uses evasions, where the most common and evolving types are timing and stalling, in response to more advanced sandboxing techniques. Evasive methods show high correlations with some malware families or packers. Despite continual technique development, there was only a modest increase in the use of evasive techniques throughout the study period. These include limiting consideration to 32-bit binaries only, the possibility of analysis tool detection by malware, and omissions regarding long-delayed behaviors’ evasions (\textbf{Galloro, Polino, Carminati, Continella, \& Zanero, 2022})\cite{Galloro}.\\

This work systematically evaluates malware detection techniques based on static analysis, dynamic analysis, and hybrid variants by training Hidden Markov Models (HMMs) on opcode and API call sequences extracted with IDA Pro (static) and Buster Sandbox Analyzer (dynamic). Using seven malware families (e.g., Zbot, ZeroAccess, Harebot, Winwebsec, Security Shield, Smart HDD) and benign Windows System32 files, the study compares four train/score settings—static-static, dynamic-dynamic, dynamic-static, and static-dynamic—under 5-fold cross-validation with ROC-AUC and Precision-Recall metrics. Results show that dynamic training and scoring deliver the strongest detection performance overall, with static-static often competitive for API calls but typically behind dynamic; simple hybrids yield mixed outcomes, where dynamic-static underperforms both pure approaches and static-dynamic performs worst. API call sequences outperform opcode sequences, likely due to opcode brittleness under obfuscation. The authors note limitations including reliance on a single sandbox/disassembly toolchain, potential labeling noise in family assignments, class imbalance affecting PR more than ROC-AUC, and uncertain generalizability beyond the tested families and Windows setting (\textbf{Damodaran, Di Troia, Corrado, Austin, \& Stamp, 2022})\cite{Damodaran}.\\

This research introduces a behavioral malware detection approach using Deep Graph Convolutional Neural Networks (GCNs) to analyze API call sequences. The study converts API call sequences into graph structures where nodes represent API calls and edges capture sequential dependencies. By applying graph convolutional layers, the model learns hierarchical representations of malware behavior patterns. Evaluated on a dataset of 42,797 malware and 1,079 goodware API call sequences extracted from Cuckoo Sandbox reports, the graph-based representation proved effective at capturing complex behavioral patterns that traditional sequence-based methods might miss. The authors note limitations including the imbalanced dataset and the computational overhead of graph construction (\textbf{Oliveira \& Sassi, 2019})\cite{Oliveira2019}.\\

Another study developed a deep learning-based sequential model for malware analysis using Windows API calls from executable files. The research utilized the Mal-API-2019 dataset containing 7,107 malware samples across eight families (Spyware, Downloader, Trojan, Worms, Adware, Dropper, Virus, Backdoor) with 342 distinct API calls. The model employed Long Short-Term Memory (LSTM) networks to capture temporal dependencies in API call sequences, treating malware detection as a sequence classification problem. The LSTM-based approach demonstrated the effectiveness of recurrent neural networks for learning from behavioral patterns in API call sequences, highlighting that sequence-based deep learning models can effectively distinguish malware families based on their runtime behavior (\textbf{Catak, Yazi, Elezaj, \& Ahmed, 2020})\cite{Catak2020}.\\

This work presents a malware classification approach using deep convolutional neural networks (CNNs) applied to malware binaries converted into grayscale images. The study transforms raw executable files into visual representations, allowing CNNs to extract spatial features and patterns characteristic of different malware families. The proposed M-CNN architecture was evaluated on two benchmark datasets: the Malimg dataset (9,339 malware images across 25 families) and the Microsoft Malware Classification Challenge dataset. The image-based approach proved highly effective for static malware analysis, as visual patterns in binary code remain relatively stable even when malware employs basic obfuscation techniques. However, the method's reliance on static features makes it vulnerable to advanced packing and encryption techniques (\textbf{Kalash, Rochan, Mohammed, Bruce, Wang, \& Iqbal, 2018})\cite{Kalash2018}.\\
\subsection{Analytical and Testing Methodologies (SAST/DAST)}
The study focuses on Android malware detection modeled through API call behavior using Markov Chains. Using a virtual environment, they ran the API calls and stored the logs. For automated input, they used the Monkey tool, which can also be useful in our concept. The approach used both static and dynamic features to analyze system behavior in different environments. They used a Random Forest Classifier to categorize the features (\textbf{Onwuzurike et al., 2018})\cite{Onwuzurike}.\\

This research focused on detecting malware in real time using a deep learning approach. It developed a dual-branch model where one branch analyzed static features, while the other looked at dynamic features. The results from both branches were combined and processed through deep neural network layers. The DNN was trained on 5.4 million PE samples to learn effective pattern representations(\textbf{Mihalache et al., 2024})\cite{RaduMihalache}.\\

Another study focused on automating DevSecOps pipelines by integrating AI. It used application testing to detect attacks and combined SAST for finding code-level vulnerabilities with DAST to catch runtime attacks like SQL injection. The system also leveraged AI-driven frameworks such as GDPR, HIPAA, OWASP, and NIST to check compliance and make sure security standards were met.This study looks at the ongoing battle between malware trying to evade detection and machine learning defenses trying to stop it. Researchers created a model that teaches malware to trick detection systems, and a defense system that trains classifiers to be tougher and more reliable against such tricks(\textbf{Abdiukov, 2024})\cite{Abdiukov}.\\

This study tackles the common problem of false alarms in detection models, which can lead to mistakes. To make the results more accurate, the researchers added an active behavioral validation step to Cloud Security Posture Management. This step simulates real attacks to test and verify AWS components like S3, EC2, and IAM settings, helping improve the overall accuracy of threat detection (\textbf{Dikshant \& Verma, 2025})\cite{Dikshant}.\\

\subsection{Adversarial and Simulation-Based Defense}
This study looks at how hackers and security systems are constantly trying to outsmart each other. The researchers build a system that shows how malware can fool machine learning–based detectors, then create another system that helps those detectors become smarter and harder to trick. This model uses API calls from PE files, simulates attacks via feature manipulation, and defends via adversarial training with security regularization. Evasion attack model achieved strong evasion with small feature changes, Secure-learning defense restored high detection accuracy, and outperformed commercial scanners. This work provides a strong foundation for building resilient, security-aware malware detection systems by successfully demonstrating that incorporating adversarial training through SecDefender can significantly strengthen malware classifiers against sophisticated evasion attacks (\textbf{Chen, Ye, \& Bourlai, 2017})\cite{Chen}.\\

This research aims to bridge the gap between AI theory and real-world cybersecurity applications by reviewing and analyzing modern cybersecurity frameworks that claim to use AI, ML, or DL for threat detection, prevention, and response. From 3,200 vendors, 38 cybersecurity frameworks were identified, focusing on IDPS, EDR, and IAM systems, and they were classified into groups — AI-powered, ML-powered, DL-powered, rules-based, or unknown/no info. Most cybersecurity frameworks are still rules-based or traditional ML-driven; few fully use deep learning, where only NVIDIA Morpheus stands out for its multi-GPU AI pipeline. Overall, the paper provides a comprehensive assessment of AI and deep learning integration within modern cybersecurity frameworks, revealing a significant gap between theoretical advances and practical deployment (\textbf{Applebaum, Miller, Strom, Korban, \& Wolf, 2016})\cite{Applebaum}.\\

Another investigation centers on black-box targeted adversarial attacks by generating more transferable adversarial samples, which proposes a new robust training framework that adjusts the decision boundaries of surrogate models to improve attack transferability. The approach is a generative adversarial framework consisting of adversarial samples from clean images, where a U-Net architecture is used, known as an attack generator, and controlled noise to smooth the surrogate model’s decision boundaries via robust training, known as a perturbator. The Discriminator (D) is robustly trained to reduce discrepancies between the surrogate and target models, improving sample transferability and Experiments are performed on ImageNet using pretrained CNNs as both surrogate and target models. Overall, the paper presents an effective and computationally efficient framework for generating transferable adversarial examples through robustly trained surrogate models, and the method significantly enhances attack success rates against black-box targets while maintaining fast generation speed (\textbf{Shi, Wang, Li, Guo, \& Liu, 2024})\cite{Shi}.\\

This work describes a Red Team/Blue Team exercise designed to test a collaborative DDoS defense system combining two mechanisms, COSSACK and D-WARD, by evaluating how well the system detects and mitigates different types of DDoS attacks, identifying vulnerabilities and weaknesses in both systems, and gathering insights for improving real-world collaborative DDoS defenses. Here blue team is the developers of COSSACK and D-WARD (defenders), the red team is the attackers from Sandia Labs, who designed and executed diverse DDoS attacks, and the white team is the Neutral coordinators from BBN Technologies, responsible for data collection and analysis. Here, nine connected edge networks are linked through a simulated core router, and realistic legitimate traffic is generated using the Skaion system. The combined COSSACK–D-WARD defense maintained high legitimate traffic and low false positives during spoofing and rolling floods, though timing-based floods caused brief disruptions. In contrast, ACK Proxy attacks effectively bypassed detection, severely degrading performance. Overall, the study provides one of the first realistic evaluations of collaborative DDoS defenses under controlled Red Team/Blue Team conditions, where the results show that COSSACK and D-WARD can effectively mitigate common flooding attacks but remain vulnerable to sophisticated proxy-based and timing attacks (\textbf{Mirkovic et al., 2008})\cite{Mirkovic}.\\



\section{Summary of Key Findings}

\setlength{\tabcolsep}{18pt}
\renewcommand{\arraystretch}{1.5}
\setlength{\arrayrulewidth}{0.5mm}

\begin{longtable}{|M{2cm}|M{3.5cm}|M{7cm}|}
\caption[Strategies and features found in papers]{Strategies and features found in papers}\label{tab:key-findings}\\
\hline
\textbf{Research Papers} & \textbf{Key Strategy} & \textbf{Feature} \\\hline
\endfirsthead

\hline
\textbf{Research Papers} & \textbf{Key Strategy} & \textbf{Feature} \\\hline
\endhead

\hline \multicolumn{2}{r}{\textit{Continued on next page}} \\\hline
\endfoot

\hline
\endlastfoot

\textcite{Modalavalasa} & AE-MLP Model (Autoencoder \& Multi-layer Perception) & AE takes the dataset, encodes insightful features, and feeds them to MLP. \\\hline
\textcite{Marandi} & Static Application Security Testing (SAST) & Uses Snyk to detect vulnerabilities during development. \\\hline
\textcite{Marandi} & Dynamic Application Security Testing (DAST) & Uses StackHawk to test and secure the running application/pipeline. \\\hline
\textcite{Ashtagi} & Software Composition Analysis (SCA) & Uses OWASP Dependency-Check to scan third-party libraries for known vulnerabilities. \\\hline
\textcite{Ashtagi}  & Container Image Scanning & Uses Trivy to scan Docker images for vulnerabilities and misconfigurations. \\\hline
\textcite{Ashtagi}  & Cloud-Native Architecture & Uses AWS VPC, EC2, CloudFront, and CloudWatch for secure, scalable, resilient deployments. \\\hline
\textcite{Ashtagi}  \textcite{Ganugapati} & Infrastructure as Code (IaC) & Uses Terraform to provision and manage AWS infrastructure. \\\hline
\textcite{Miikkulainen} \textcite{Harvey} & Evolving Neural Networks (CNE/NEAT, indirect encoding) & Evolves topology and/or weights; robust in non-stationary tasks; indirect encoding scales to large networks. \\\hline
\textcite{Kotyan} & Robust NAS (R-NAS) & NAS jointly optimizes accuracy and adversarial robustness; improves resistance to black-box attacks. \\\hline
\textcite{Tsukamoto} & Adaptive neural model with ``neuro-agents'' & Modular, hierarchical neuro-objects evolve structure and learning parameters in parallel. \\\hline
\textcite{Christodorescu} & Semantics-aware instruction-template detection & Architecture-independent IR + symbolic execution/theorem proving; fewer signatures, low false positives. \\\hline
\textcite{Kolosnjaji} & ConvNet-LSTM on system-call sequences & Dynamic, sequence-based malware classification from sandboxed API calls. \\\hline
\textcite{Galloro} & Evasion tactics landscape & Widespread sandbox/debugger evasion (e.g., timing, stalling); need behaviorally robust defenses. \\\hline
\textcite{Damodaran} & Static vs. dynamic vs. hybrid detection & Dynamic analysis generally outperforms static; API-call features often beat opcode features. \\\hline
\textcite{Oliveira2019} & Graph CNN on API call sequences & Converts API sequences to graphs; GCN learns behavioral patterns from imbalanced dataset. \\\hline
\textcite{Catak2020} & LSTM on API call sequences (Mal-API-2019) & Sequential deep learning on malware samples; effective for temporal behavior patterns. \\\hline
\textcite{Kalash2018} & M-CNN on malware images (Malimg) & Converts binaries to grayscale images; CNN extracts visual patterns for classification. \\\hline
\textcite{Chen} & Adversarial training (SecDefender) & Security-aware training increases robustness to feature-manipulation/evasion attacks. \\\hline
\textcite{Applebaum} & State of AI in cyber products & Many frameworks remain rules/ML-based; limited DL adoption; motivates research gap. \\\hline
\textcite{Shi} & Transferable adversarial examples (black-box) & Robustly trained surrogate models (generator/perturbator/discriminator) yield transferable attacks. \\\hline
\textcite{Mirkovic} & Red/Blue-team evaluation (COSSACK \& D-WARD) & Effective vs. spoofing/rolling floods; gaps vs. ACK-proxy/timing attacks. \\\hline
\end{longtable}
